\section{Conclusions and Future Work}

In this work, we demonstrated an open-source workflow for high-throughput microstructure stability analysis. Firstly, if a suitable CALPHAD model is not available to the modeler, it can be created through PhaseForge. The chemical Hessian is then determined analytically and efficiently via Jansson derivatives. Coherent strain contributions can be added through an analytically differentiable bond model and MLIP estimates for elastic constants. Lastly, following a clever change of basis, gauge-invariant eigenvalues and eigenvectors of the chemical and coherent spinodal can be investigated at specific compositions as functions of temperature or visualized at a given temperature over a high-dimensional composition space through affine projections. This workflow provides many valuable extensions to the foundational stability analysis framework produced by Koneru, Kadirvel, and Wang \cite{koneru2022high}, to include 1) open-source access, 2) whitebox CALPHAD model knowledge, 3) improved efficiency and robustness of Gibbs energy derivatives, 4) modeling coherent strain contributions, and 5) visualization in large composition spaces via one two-dimensional plot.

However, there are still plenty of future directions to further extend this workflow. A model for homogeneous elastic Gibbs energy contributions could be incorporated in addition to contributions from coherent strain due to composition fluctuations. MLIPs could also be used to create bond models for other crystal structures, and coherent strain elastic contributions could be modeled for materials that do not behave isotropically. Lastly, implementation of second-order Jansson derivatives would provide analytic calculation of third-order derivatives of the Gibbs energy with respect to composition, enabling the analytic calculation of critical points.
